
en notant que \(i_x:T\to T\times_k C\) est un morphisme \emph{parfait}, car déduit par le changement de base plat \(T\times_k C\to C\) du morphisme \(j_x:x\to C\), qui est parfait grâce à l'hypothèse \(C\) régulier. On trouve ainsi un groupe quotient
\begin{equation}\tag{7.11.3}
  K_\bullet^{\mathrm{alg}}(T)
\end{equation}
de \(K_\bullet(T)\). Ce dernier est d'ailleurs de façon naturelle un module sur \(K^\bullet_{\mathrm{alg}}(T)\), car on constate aussitôt sur les définitions que le produit d'un éléments de \(K^\bullet(T)\) par un élément de \(K_\bullet(T)\) est algébriquement équivalent à zéro si l'un des deux facteurs l'est, en utilisant (IV 2.12.2) pour \(i_x\).

On définit (quand on travaille avec le \(K^\bullet\) naïf) l'idéal de \(\Gr^\bullet(T)\) des éléments algébriquement équivalents à zéro comme étant l'idéal engendré par les \(c^i(x)\), où \(i\geq1\) et \(x\in K^\bullet(T)\) est algébriquement équivalent à zéro; l'anneau gradué quotient
\begin{equation}\tag{7.11.4}
  \Gr^\bullet_{\mathrm{alg}}(T)
\end{equation}
n'est alors autre que l'anneau de Chern associé au \(\lambda\)-anneau \(K^\bullet_{\mathrm{alg}}(T)\).

Enfin on définit le sous-groupe gradué de \(\Gr_\bullet(T)\) des éléments algébriquement équivalents à zéro (\(T\) étant supposé localement noethérien), comme celui dont la composante de dimension \(i\) est l'image dans \(\Gr_i(T)\) du sous-groupe de \(\Fil_iK_\bullet(T)\) formé des éléments algébriquement équivalents à zéro dans \(K_\bullet(T)\). On trouve aussitôt que le groupe gradué quotient
\begin{equation}\tag{7.11.5}
  \Gr_\bullet^{\mathrm{alg}}(T)
\end{equation}
est un module sur l'anneau \(\Gr^\bullet_{\mathrm{alg}}(T)\), par passage au quotient à partir de la structure de module de \(\Gr_\bullet(T)\) sur \(\Gr^\bullet(T)\) (3.1).

\textbf{7.11.6.} Les définitions qui précèdent garderaient un sens si on ne suppose pas \(k\) algébriquement clos, en prenant ci-dessus \(C\) \emph{géométriquement} connexe sur \(k\). Il n'est pas clair cependant que ces définitions seraient bien utiles; notamment, il ne semble pas que dans le cas du groupe \(\Pic(T)\), on
